%File: Wheelchair_accesible_urban_navigation.tex
\documentclass[letterpaper]{article} % DO NOT CHANGE THIS
\usepackage{aaai2026}  % DO NOT CHANGE THIS
\usepackage{times}  % DO NOT CHANGE THIS
\usepackage{helvet}  % DO NOT CHANGE THIS
\usepackage{courier}  % DO NOT CHANGE THIS
\usepackage[hyphens]{url}  % DO NOT CHANGE THIS
\usepackage{graphicx} % DO NOT CHANGE THIS
\urlstyle{rm} % DO NOT CHANGE THIS
\def\UrlFont{\rm}  % DO NOT CHANGE THIS
\usepackage{natbib}  % DO NOT CHANGE THIS AND DO NOT ADD ANY OPTIONS TO IT
\usepackage{caption} % DO NOT CHANGE THIS AND DO NOT ADD ANY OPTIONS TO IT
\frenchspacing  % DO NOT CHANGE THIS
\setlength{\pdfpagewidth}{8.5in} % DO NOT CHANGE THIS
\setlength{\pdfpageheight}{11in} % DO NOT CHANGE THIS

% Recommended packages for algorithms
\usepackage{algorithm}
\usepackage{algorithmic}

% Recommended packages for listings (code blocks)
\usepackage{newfloat}
\usepackage{listings}
\DeclareCaptionStyle{ruled}{labelfont=normalfont,labelsep=colon,strut=off}
\lstset{%
	basicstyle={\footnotesize\ttfamily},
	numbers=left,numberstyle=\footnotesize,xleftmargin=2em,
	aboveskip=0pt,belowskip=0pt,
	showstringspaces=false,tabsize=2,breaklines=true}
\floatstyle{ruled}
\newfloat{listing}{tb}{lst}{}
\floatname{listing}{Listing}

\pdfinfo{
/TemplateVersion (2026.1)
}



\nocopyright

\setcounter{secnumdepth}{0}

\usepackage{amsmath}


% TITLE AND AUTHORS

\title{AI Techniques for Wheelchair-Accessible Urban Navigation}

\author{
    Fatima Surname\equalcontrib,
    Francesca Surname\equalcontrib,
    Lara Surname\equalcontrib,
    Matteo Surname\equalcontrib
}
\affiliations{
    Université libre de Bruxelles (ULB)\\
    INFO-H410 --- Techniques of Artificial Intelligence
}

% BEGIN DOCUMENT
\begin{document}

\maketitle

\begin{abstract}
% (1) Problema y motivación.
% (2) Aportación: comparamos Search, CSP y MDP sobre el mismo problema.
% (3) Método: implementación + métricas de evaluación.
% (4) Resultado clave / cuándo cada enfoque es preferible.
\end{abstract}

\begin{links}
    \link{Code}{https://github.com/francescagrasso02/TechniquesOfAI}
\end{links}

\section{Introduction}

Currently, route planning applications are mainly based on static criteria such as minimum distance or minimum estimated walking time. However, for wheelchair users, these parameters are not sufficient to be able to find a route that they can take independently. Factors such as narrow sidewalks, steep ramps or inaccessible kerbs may make the route impossible for people with mobility impairments \cite{darko2022}. Recent studies on urban mobility in Brussels highlight problems related to inaccessible kerbs, discontinuities in the tactile pavement and obstacles present on sidewalks \cite{dobruszkes2025}.

\section{Problem Formulation}

This work proposes a route planning system that focuses on accessibility for wheelchair users in the city of Brussels. Using real OpenStreetMap data, this system models the urban pedestrian network as a graph, taking into account crucial accessibility factors such as slope, sidewalk width, surface type and kerb accessibility \cite{kaselouris2025}.

Three distinct artificial intelligence approaches are implemented to solve this problem: heuristic search algorithms, such as Dijkstra and A*, Constraint Satisfaction Problems (CSP) and Markov Decision Processes (MDP). These three approaches, which address the same problem from different perspectives, will be compared experimentally and qualitatively to analyze their strengths, and limitations.

\subsection{Formal Representation}

The pedestrian network is modeled as a weighted graph:
\[
G = (V,E)
\]

where:
\begin{itemize}

    \item \(V\) represents navigable pedestrian nodes extracted from OpenStreetMap
    
    \item \(E\) represents sidewalk or street segments connecting adjacent nodes
    
\end{itemize}

Each edge \(e \in E\) contains accessibility-related attributes:
\[
e = (\text{length}, \text{slope}, \text{width}, \text{kerb}, \text{wheelchair}, \text{surface})
\]

including:
\begin{itemize}
    \item \textbf{length}: street length in metres
    \item \textbf{slope}: gradient percentage
    \item \textbf{width}: sidewalk width in metres
    \item \textbf{kerb}: kerb accessibility type
    \item \textbf{wheelchair}: wheelchair accessibility tag (\textit{yes, no, limited, unknown})
    \item \textbf{surface}: surface type
\end{itemize}

The objective is to compute an accessible path:

\[
P = (v_1,v_2,\dots,v_n)
\]

between an origin node \(v_1\) and a destination node \(v_n\), while minimizing traversal cost, which is primarily based on path length:

\[
\min \sum_{e \in P} c(e)
\]

The routing process is subject to several accessibility constraints related to pedestrian mobility, including sidewalk slope, minimum width requirements, kerb accessibility, and wheelchair accessibility tags. The system applies strict constraints by default, but switches to relaxed constraints when no feasible accessible route can be found.

\subsection{Suitability for Multiple AI Techniques}
This problem allows the application of search algorithms such as Dijkstra and A*, since the pedestrian network can be represented as a weighted graph for finding minimum-cost accessible routes. Constraint Satisfaction Problems can also be applied to ensure that the route satisfies specific accessibility constraints. Finally, Markov Decision Processes can be used to model uncertainty in accessibility conditions due to the presence of incomplete accessibility information in OpenStreetMap data.


\section{Methodological Approaches}

\subsection{Approach A: Search Algorithms}

\subsection{Approach B: Constraint Satisfaction Problems}

\subsection{Approach C: Markov Decision Processes}

\section{Experimental Evaluation}
% métricas, instancias, tabla comparativa, figuras.

\section{Qualitative Comparison}

\section{Conclusion}

\bibliography{references}

\end{document}